\documentclass[10pt,a4paper]{amsart}
\usepackage{amsmath,amssymb,amsthm,graphicx}
\usepackage{verbatim}
\newenvironment{reference}{\comment}{\endcomment}
%\newenvironment{reference}{}{}
\newenvironment{slogan}{\comment}{\endcomment}
\newenvironment{history}{\comment}{\endcomment}

% For commutative diagrams we use Xy-pic
\usepackage[all]{xy}

% We use 2cell for 2-commutative diagrams.
\xyoption{2cell}
% Theorem environments.
%
\theoremstyle{plain}
\newtheorem{theorem}[subsection]{Theorem}
\newtheorem{proposition}[subsection]{Proposition}
\newtheorem{lemma}[subsection]{Lemma}

\theoremstyle{definition}
\newtheorem{definition}[subsection]{Definition}
\newtheorem{example}[subsection]{Example}
\newtheorem{exercise}[subsection]{Exercise}
\newtheorem{situation}[subsection]{Situation}

\theoremstyle{remark}
\newtheorem{remark}[subsection]{Remark}
\newtheorem{remarks}[subsection]{Remarks}

\numberwithin{equation}{subsection}

% Macros
%
\def\lim{\mathop{\mathrm{lim}}\nolimits}
\def\colim{\mathop{\mathrm{colim}}\nolimits}
\def\Spec{\mathop{\mathrm{Spec}}}
\def\Hom{\mathop{\mathrm{Hom}}\nolimits}
\def\Ext{\mathop{\mathrm{Ext}}\nolimits}
\def\SheafHom{\mathop{\mathcal{H}\!\mathit{om}}\nolimits}
\def\SheafExt{\mathop{\mathcal{E}\!\mathit{xt}}\nolimits}
\def\Sch{\mathit{Sch}}
\def\Mor{\mathop{\mathrm{Mor}}\nolimits}
\def\Ob{\mathop{\mathrm{Ob}}\nolimits}
\def\Sh{\mathop{\mathit{Sh}}\nolimits}
\def\NL{\mathop{N\!L}\nolimits}
\def\CH{\mathop{\mathrm{CH}}\nolimits}
\def\proetale{{pro\text{-}\acute{e}tale}}
\def\etale{{\acute{e}tale}}
\def\QCoh{\mathit{QCoh}}
\def\Ker{\mathop{\mathrm{Ker}}}
\def\Im{\mathop{\mathrm{Im}}}
\def\Coker{\mathop{\mathrm{Coker}}}
\def\Coim{\mathop{\mathrm{Coim}}}

% Boxtimes
%
\DeclareMathSymbol{\boxtimes}{\mathbin}{AMSa}{"02}

%
% Macros for moduli stacks/spaces
%
\def\QCohstack{\mathcal{QC}\!\mathit{oh}}
\def\Cohstack{\mathcal{C}\!\mathit{oh}}
\def\Spacesstack{\mathcal{S}\!\mathit{paces}}
\def\Quotfunctor{\mathrm{Quot}}
\def\Hilbfunctor{\mathrm{Hilb}}
\def\Curvesstack{\mathcal{C}\!\mathit{urves}}
\def\Polarizedstack{\mathcal{P}\!\mathit{olarized}}
\def\Complexesstack{\mathcal{C}\!\mathit{omplexes}}
% \Pic is the operator that assigns to X its picard group, usage \Pic(X)
% \Picardstack_{X/B} denotes the Picard stack of X over B
% \Picardfunctor_{X/B} denotes the Picard functor of X over B
\def\Pic{\mathop{\mathrm{Pic}}\nolimits}
\def\Picardstack{\mathcal{P}\!\mathit{ic}}
\def\Picardfunctor{\mathrm{Pic}}
\def\Deformationcategory{\mathcal{D}\!\mathit{ef}}
\usepackage{cleveref}
\crefname{originalreference}{original reference}{original references}
\Crefname{originalreference}{Original reference}{Original references}
\makeatletter
\newlabel{stacks-morphisms-lemma-composition-smooth}{{075V}{source}{}{sourcepointer-1}{}}
\newlabel{stacks-morphisms-lemma-composition-smooth@cref}{{[originalreference][0][]075V}{[1][0][]source}{}{}{}}
\makeatother
\crefname{originalreference}{original reference}{original references}
\Crefname{originalreference}{Original reference}{Original references}
\makeatletter
\newlabel{stacks-sheaves-lemma-pullback-lqc}{{06WL}{source}{}{sourcepointer-2}{}}
\newlabel{stacks-sheaves-lemma-pullback-lqc@cref}{{[originalreference][0][]06WL}{[1][0][]source}{}{}{}}
\makeatother
\crefname{originalreference}{original reference}{original references}
\Crefname{originalreference}{Original reference}{Original references}
\makeatletter
\newlabel{stacks-sheaves-lemma-surjective-flat-locally-finite-presentation}{{06XH}{source}{}{sourcepointer-3}{}}
\newlabel{stacks-sheaves-lemma-surjective-flat-locally-finite-presentation@cref}{{[originalreference][0][]06XH}{[1][0][]source}{}{}{}}
\makeatother
\crefname{originalreference}{original reference}{original references}
\Crefname{originalreference}{Original reference}{Original references}
\makeatletter
\newlabel{stacks-sheaves-lemma-comparison}{{075D}{source}{}{sourcepointer-4}{}}
\newlabel{stacks-sheaves-lemma-comparison@cref}{{[originalreference][0][]075D}{[1][0][]source}{}{}{}}
\makeatother
\crefname{originalreference}{original reference}{original references}
\Crefname{originalreference}{Original reference}{Original references}
\makeatletter
\newlabel{spaces-properties-lemma-characterize-quasi-coherent}{{03M0}{source}{}{sourcepointer-5}{}}
\newlabel{spaces-properties-lemma-characterize-quasi-coherent@cref}{{[originalreference][0][]03M0}{[1][0][]source}{}{}{}}
\makeatother
\title{Local quasi-coherence descends along smooth coverings}
\author{The Stacks Project Authors}
\date{stacks-project:stacks-cohomology.tex:local-snapshot-20261012; GNU FDL 1.2 or later}
\begin{document}
\maketitle
\noindent\textit{Editorial note.} All delivered mathematical text is the author's own, copied verbatim from the pinned Stacks Project chapter file stacks-cohomology.tex; no statement or proof step was rewritten, reordered or removed.\par\medskip
\noindent\textit{Editorial note.} The counted span cites the lemmas recorded in assembly.external\_references with canonical tags as navigation pointers.\par\medskip
\noindent\textit{Allowed standard prerequisites (editorial).}\begin{itemize}
\item Surjective flat lfp lifting. Sheaves on Stacks, Lemma stacks-sheaves-lemma-surjective-flat-locally-finite-presentation. Etale covering with lifts.
\item Comparison of restrictions. Sheaves on Stacks, Lemma stacks-sheaves-lemma-comparison. x\_i\textasciicircum{}*F as restriction of f\textasciicircum{}*F.
\item Characterization of qc on algebraic spaces. Properties of Spaces, Lemma spaces-properties-lemma-characterize-quasi-coherent. Final step.
\end{itemize}\medskip
\ifdefined\hypertarget\hypertarget{sourcepointer-1}{}\fi
\noindent\textit{Source reference.} Original reference 075V: sources/stacks-chow/tags.txt, source lines 9144--9144. Cross-chapter standard prerequisite cited by the author inside the counted span; it is used as a pointer and its own proof is not part of this item.\par\medskip
\ifdefined\hypertarget\hypertarget{sourcepointer-2}{}\fi
\noindent\textit{Source reference.} Original reference 06WL: sources/stacks-chow/tags.txt, source lines 8820--8820. Cross-chapter standard prerequisite cited by the author inside the counted span; it is used as a pointer and its own proof is not part of this item.\par\medskip
\ifdefined\hypertarget\hypertarget{sourcepointer-3}{}\fi
\noindent\textit{Source reference.} Original reference 06XH: sources/stacks-chow/tags.txt, source lines 8851--8851. Cross-chapter standard prerequisite cited by the author inside the counted span; it is used as a pointer and its own proof is not part of this item.\par\medskip
\ifdefined\hypertarget\hypertarget{sourcepointer-4}{}\fi
\noindent\textit{Source reference.} Original reference 075D: sources/stacks-chow/tags.txt, source lines 9127--9127. Cross-chapter standard prerequisite cited by the author inside the counted span; it is used as a pointer and its own proof is not part of this item.\par\medskip
\ifdefined\hypertarget\hypertarget{sourcepointer-5}{}\fi
\noindent\textit{Source reference.} Original reference 03M0: sources/stacks-chow/tags.txt, source lines 4732--4732. Cross-chapter standard prerequisite cited by the author inside the counted span; it is used as a pointer and its own proof is not part of this item.\par\medskip
\begin{lemma}
\label{lemma-check-lqc-on-etale-covering}
Let $\mathcal{X}$ be an algebraic stack. Let
$f_j : \mathcal{X}_j \to \mathcal{X}$ be a family of smooth
morphisms of algebraic stacks with
$|\mathcal{X}| =\bigcup |f_j|(|\mathcal{X}_j|)$.
Let $\mathcal{F}$ be a sheaf of $\mathcal{O}_\mathcal{X}$-modules
on $\mathcal{X}_\etale$. If each $f_j^{-1}\mathcal{F}$
is locally quasi-coherent, then so is $\mathcal{F}$.
\end{lemma}

\begin{proof}
We may replace each of the algebraic stacks $\mathcal{X}_j$ by
a scheme $U_j$ (using that any algebraic stack has a smooth covering by
a scheme and that compositions of smooth morphisms are smooth, see
Morphisms of Stacks, Lemma \ref{stacks-morphisms-lemma-composition-smooth}).
The pullback of $\mathcal{F}$ to $(\Sch/U_j)_\etale$ is still
locally quasi-coherent, see
Sheaves on Stacks, Lemma \ref{stacks-sheaves-lemma-pullback-lqc}.
Then $f = \coprod f_j : U = \coprod U_j \to \mathcal{X}$ is a surjective
smooth morphism. Let $x$ be an object of $\mathcal{X}$. By
Sheaves on Stacks, Lemma
\ref{stacks-sheaves-lemma-surjective-flat-locally-finite-presentation}
there exists an \'etale covering $\{x_i \to x\}_{i \in I}$
such that each $x_i$ lifts to an object $u_i$ of $(\Sch/U)_\etale$.
This just means that $x$, $x_i$ live over schemes $V$, $V_i$, that
$\{V_i \to V\}$ is an \'etale covering, and that $x_i$ comes from
a morphism $u_i : V_i \to U$. The restriction
$x_i^*\mathcal{F}|_{V_{i, \etale}}$ is equal to the restriction
of $f^*\mathcal{F}$ to $V_{i, \etale}$, see
Sheaves on Stacks, Lemma \ref{stacks-sheaves-lemma-comparison}.
Hence $x^*\mathcal{F}|_{V_\etale}$
is a sheaf on the small \'etale site of $V$ which is quasi-coherent
when restricted to $V_{i, \etale}$ for each $i$.
This implies that it is quasi-coherent (as desired), for example by
Properties of Spaces, Lemma
\ref{spaces-properties-lemma-characterize-quasi-coherent}.
\end{proof}

\end{document}
